\begin{proposition} \label{proposition:natural_numbers_multiplication_associative_commutative_distributive}
Let $\bbN$ be the set of \CrefAndHyperrefIfExist{definition:natural_numbers}{natural numbers}. For all $n, m, k \in \bbN$:
\begin{enumerate}
    \item $(n \times m) \times k = n \times (m \times k)$ (associativity),
    \item $n \times m = m \times n$ (commutativity),
    \item $n \times (m + k) = n \times m + n \times k$ (distributivity).
\end{enumerate}
\CrefIfExists{proposition:natural_numbers_addition_existence}\CrefIfExists{proposition:natural_numbers_multiplication_existence}
Moreover, $1$ is the multiplicative identity: $n \times 1 = n$ for all $n \in \bbN$.
\end{proposition}

\begin{proof}

        We first show the distributivity property by fixing $n,m \in \bbN$ and inducting on $k$.  
        \begin{itemize}
            \item For the base case, i.e. $k = 1$, we have 
            $$n \times (m+1) = n \times s(m) = n \times m + n = n \times m + n \times 1$$
            by \Cref{proposition:natural_numbers_addition_existence} and \Cref{proposition:natural_numbers_multiplication_existence}.

            \item Suppose inductively that $l \in \bbN$ satisfies
            $$n \times (m+l) = n \times m + n \times l.$$
            We need to show that 
            $$n \times (m+s(l)) = n \times m + n \times s(l).$$
            The left hand side equals
            $$n \times s(m+l) = n \times (m+l) + n = (n \times m + n \times l) + n$$
            by \Cref{proposition:natural_numbers_addition_existence} and the inductive hypothesis. By the associtiativy property in \Cref{proposition:natural_numbers_addition_associative_commutative}, this equals 
            $$(n \times m) + (n \times l + n)$$
            which equals
            $$(n \times m) + (n \times s(l))$$
            as desired.
        \end{itemize}


        We show that $\times$ is associative by fixing $n, m$ and inducting on $k$.
        \begin{itemize}
            \item For the base case, i.e. $k = 1$, by \Cref{proposition:natural_numbers_multiplication_existence}, we have
            $$(n \times m) \times 1 = n \times m$$
            and
            $$n \times (m \times 1) = m \times m$$
            Therefore, 
            $$(n \times m) \times 1 = n \times (m \times 1).$$

            \item Suppose inductively that $l \in \bbN$ is an integer such that 
            $$(n \times m) \times l = n \times (m \times l);$$
            we aim to show that 
            $$(n \times m) \times s(l) = n \times (m \times s(l)).$$
            \CrefIfExists{definition:natural_numbers}
            By \Cref{proposition:natural_numbers_multiplication_existence}, 
            $$(n \times m) \times s(l) = (n \times m) \times l + (n \times m),$$
            which by the inductive hypothesis equals 
            $$n \times (m \times l) + n \times m.$$
            In turn, by distributivity, this equals 
            $$n \times (m \times l + m)$$
            and once again, by \Cref{proposition:natural_numbers_multiplication_existence}, this equals 
            $$n \times (m \times s(l)).$$
            We have thus shown the desired equality.
        \end{itemize}

        We show that $\times$ is commutative by inducting on $n$ for fixed $m$.

        \begin{itemize}
            \item For the base case, i.e. $n = 1$, $1 \times m = m$ by \Cref{lemma:natural_number_one_left_identity_for_multiplication}, and $m \times 1 = m$ by \CrefIfExists{proposition:natural_numbers_multiplication_existence}. Hence $1 \times m = m \times 1$.

            \item Suppose that $k \in \bbN$ satisfies $k \times m = m \times k$. By \Cref{lemma:natural_number_successor_left_multiplication_recursion}, the inductive hypothesis, and \CrefIfExists{proposition:natural_numbers_multiplication_existence}, we have
            $$s(k) \times m = k \times m + m = m \times k + m = m \times s(k).$$
        \end{itemize}

\end{proof}
